%=============================================================================
%  Programme brief for the faculty mentor.
%  Plain language, no competition framing. Build:  pdflatex mentor-brief
%=============================================================================
\documentclass[11pt,a4paper]{article}
\usepackage[T1]{fontenc}
\usepackage{mathptmx}
\usepackage[margin=24mm]{geometry}
\usepackage{booktabs}
\usepackage{array}
\usepackage{microtype}
\usepackage{titlesec}
\usepackage[hidelinks]{hyperref}
\usepackage{fancyhdr}

\newcolumntype{L}[1]{>{\raggedright\arraybackslash}p{#1}}
\newcommand{\rs}[1]{\mbox{INR~#1}}
\titleformat{\section}{\large\bfseries}{\thesection.}{0.6em}{}
\setlength{\parskip}{0.55em}
\setlength{\parindent}{0pt}
\pagestyle{fancy}\fancyhf{}
\fancyhead[L]{\small Flood search and rescue --- programme brief}
\fancyhead[R]{\small\thepage}
\renewcommand{\headrulewidth}{0.4pt}
\renewcommand{\arraystretch}{1.2}

\begin{document}

\begin{center}
{\LARGE\bfseries Finding People in Floodwater from the Air}\\[3mm]
{\large A three-drone search system, the flaw we found in it,\\
and what it will take to finish}\\[5mm]
Swastik Kumar\\
Department of Electronics \& Communication Engineering\\
Thapar Institute of Engineering \& Technology\\[2mm]
\today
\end{center}
\thispagestyle{empty}
\vspace{2mm}\hrule\vspace{5mm}

%-----------------------------------------------------------------------------
\section{Executive summary}

Floods displace millions of people in India every year. The slow part of a
rescue is usually not reaching survivors --- it is \emph{finding} them.

We are building three one-metre, 6.4\,kg multirotor aircraft that fly a search
pattern over flooded ground, recognise people from the air, report their
coordinates, and drop a small relief kit to each. Everything has to be decided
on a computer small enough to fly, without any phone or internet connection.

Most of the design is settled and its arithmetic is written down and
reproducible. But while checking it we found a real flaw: \textbf{our camera
sees a survivor clearly, and our own software throws that detail away before
the detector looks at it.} We have a fix, and we can test it on free public data
before committing to anything.

This document explains the flaw, the fix, and what the whole programme needs ---
\textbf{\rs{7{,}34{,}379} across seven small steps}, each released only when
the previous one has produced a result. The first decision is INR~31,371.

%-----------------------------------------------------------------------------
\section{Problem}

\textbf{We are discarding resolution we already paid for.}

A person in floodwater shows only their head and shoulders --- roughly 40\,cm
across. Our camera, flying at 60\,m, resolves that across about 26 pixels. But
detection models take a fixed, small input, so the picture is shrunk to fit.
After shrinking, the person occupies about \textbf{13 pixels}.

At that size, published results for detecting tiny people report accuracy
between 29\,\% and 61\,\%. That is not good enough when a missed person is a
missed rescue. And the obvious response --- use a smaller, faster model ---
makes it \emph{worse}, because smaller models take smaller inputs.

Three further things make floods harder than they first appear:

\begin{itemize}\itemsep2pt
\item \textbf{No suitable training data exists.} Flood datasets label buildings
and water damage, not people. Datasets of people in water are open sea: clear
blue water, no debris, no rooftops.
\item \textbf{Thermal cameras do not rescue us.} In water, skin cools toward
water temperature and the heat difference largely disappears. Thermal helps on
rooftops and fails in exactly the case we care most about.
\item \textbf{We were sizing for the wrong target.} Our early work assumed a
person lying flat, 1.7\,m long, rather than a person in water, 0.4\,m across.
That made the problem look four times easier than it is. We found the mistake
ourselves --- and it is precisely why we now want to measure rather than assume.
\end{itemize}

%-----------------------------------------------------------------------------
\section{Solution}

Three changes, in order of how much they help.

\textbf{1. Fly lower.} Coming down from 60\,m to 30\,m more than triples how
large a person appears. It costs a little flying time, and we have a great deal
to spare --- the planned mission takes under eight minutes.

\textbf{2. Stop shrinking the picture.} Instead, cut the full-resolution image
into overlapping tiles and examine each one properly. This is well established
practice. The catch is that it costs four times the computing power.

\textbf{3. Make step 2 affordable with a cheap first look.} Most of a flooded
area is empty water. A small, fast model checks each tile and asks only
\emph{``could there be anything here?''} The expensive detector then runs only
where the answer is yes. We calculate that about \textbf{87\,\% of tiles are
empty}, and we need to skip \textbf{80\,\%} to break even --- so the arithmetic
closes.

Steps 1 and 2 together take a person from \textbf{13 pixels to 52} ---
\textbf{sixteen times the area} --- for the same computing cost.

\textbf{The honest risk.} The cheap first look might discard a tile that had a
person in it. That failure is silent from the air and cannot be undone later.
This is exactly why we intend to test it before adopting it, and why our pass
mark assumes the pessimistic case: that such mistakes repeat from frame to frame
rather than cancelling out.

%-----------------------------------------------------------------------------
\section{Deliverables}

\begin{enumerate}\itemsep3pt
\item \textbf{Three working aircraft}, flown as a coordinated group under one
operator, with no phone or internet dependency.
\item \textbf{A clear yes or no on the detection method}, judged against a
threshold we wrote down \emph{before} running anything: skip at least 80.5\,\%
of tiles while still finding at least 90\,\% of people.
\item \textbf{A measured detection figure} to replace the calculated one we rely
on today. It would be the first real number we have.
\item \textbf{An evidence-based answer on flight altitude}, which our own design
currently leaves unsettled between 60\,m and 40\,m.
\item \textbf{A possible publication.} Three research areas --- image tiling,
aerial person detection, and low-power onboard processing --- have not been
combined for flood conditions. That is a narrow, defensible contribution, and it
is testable rather than merely argued.
\item \textbf{Code, data and results published either way.} A negative result
still settles an open question and costs the same to obtain.
\end{enumerate}

%-----------------------------------------------------------------------------
\section{Resources needed}

Money is not needed in one decision, or even in four. The programme splits into
\textbf{seven small steps}, and the order is engineering rather than accounting:
prove the motor before buying twelve of them, fly one aircraft before building
three.

Each step is released only when the previous one has produced a result, so the
department is never committed to more than one step at a time --- and the first
decision is \rs{31{,}371}, not \rs{7{,}34{,}379}.

All prices are current Indian retail. Totals include customs duty, GST and 15\,\%
contingency where they apply.

\begin{center}
\begin{tabular}{@{}clrL{6.0cm}@{}}
\toprule
\textbf{\#} & \textbf{Step} & \textbf{INR} & \textbf{Released only when} \\
\midrule
1 & Prove the propulsion & 31,371 & On approval \\
2 & Aircraft one: avionics and sensing & 1,05,172 & \textbf{Measured thrust $\geq$ 3.18\,kgf per motor} \\
3 & Aircraft one: airframe and power & 89,192 & Avionics powered and communicating on the bench \\
4 & Fly aircraft one & 42,574 & Airframe assembled and weighed \\
5 & Ground station and registration & 66,992 & First flight completed \\
6 & Aircraft two & 1,99,539 & Aircraft one has flown a full mission \\
7 & Aircraft three & 1,99,539 & Two aircraft flying with verified separation \\
\midrule
& \textbf{Total} & \textbf{7,34,379} & \\
\bottomrule
\end{tabular}
\end{center}

\subsection*{Step 1 --- Prove the propulsion \quad \textbf{\rs{31{,}371}}}

The smallest step, and the one that decides everything after it. Our motors are
sold without a published thrust curve, so we do not yet know that the aircraft
can lift itself. Rather than buy twelve on faith, we buy \emph{one} and measure.

\begin{center}\small
\begin{tabular}{@{}L{3.2cm}L{4.6cm}rrL{4.0cm}@{}}
\toprule
\textbf{Item} & \textbf{Model} & \textbf{Qty} & \textbf{Total} & \textbf{Why} \\
\midrule
Calibrated scale & Bench scale, 30\,kg $\times$ 1\,g, certified & 1 & 12,000 & Mass is a hard limit; we cannot argue with a scale we do not own \\
Thrust stand & 20\,kg load cell + HX711 + power meter, frame built in-house & 1 & 8,000 & The instrument that answers the question \\
Motor & Tarot TL96020, 5008 / 340\,KV & \textbf{1} & 4,500 & One, not twelve. Must lift 3.18\,kgf \\
\bottomrule
\end{tabular}
\end{center}

\textbf{If this step fails, the programme stops here having spent
\rs{31{,}371}} --- and both instruments remain useful departmental equipment.

\subsection*{Step 2 --- Aircraft one: avionics and sensing \quad \textbf{\rs{1{,}05{,}172}}}

\begin{center}\small
\begin{tabular}{@{}L{3.2cm}L{4.6cm}rrL{4.0cm}@{}}
\toprule
\textbf{Item} & \textbf{Model} & \textbf{Qty} & \textbf{Total} & \textbf{Why} \\
\midrule
Flight controller & Holybro Pixhawk 6C Mini & 1 & 22,600 & Dual sensors, industry standard, open-source autopilot \\
GNSS receiver & Centimetre-accurate (RTK) & 1 & 18,000 & Position accuracy is the largest single error in the system \\
Camera + lens & Arducam IMX477 + 6\,mm, \textbf{fixed focus} & 1 & 11,100 & Fixed focus is correct at every altitude we fly \\
AI accelerator & Raspberry Pi AI HAT+ (Hailo-8) & 1 & 11,000 & Runs the detector on board \\
Radio + storage & ExpressLRS 2.4\,GHz & 1 & 8,500 & One link carries control \emph{and} telemetry \\
Onboard computer & Raspberry Pi 5, 8\,GB & 1 & 8,000 & The host everything runs on \\
Mesh radio & 5.8\,GHz + antennas & 1 & 6,000 & Video and data between aircraft \\
\bottomrule
\end{tabular}
\end{center}

\subsection*{Step 3 --- Aircraft one: airframe and power \quad \textbf{\rs{89{,}192}}}

\begin{center}\small
\begin{tabular}{@{}L{3.2cm}L{4.6cm}rrL{4.0cm}@{}}
\toprule
\textbf{Item} & \textbf{Model} & \textbf{Qty} & \textbf{Total} & \textbf{Why} \\
\midrule
Motors & Tarot TL96020, remaining three & 3 & 13,500 & Bought only after step 1 proved the first one \\
Airframe & 25\,$\times$\,23\,mm carbon tube (Kineco Kaman) + machined clamps & 1 & 13,000 & Fabricated in the institute machine shop \\
Battery cells & Molicel P45B 21700, 4500\,mAh & 18 & 12,600 & Peak draw is 38.3\,A per cell; a 40\,A cell has no margin \\
Pack assembly & 6S 60\,A BMS + power board + regulator & 1 & 8,500 & 18 cells must be welded, fused and balanced \\
Speed controllers & 50--60\,A continuous, 6S & 4 & 7,200 & Peak is 29\,A per motor \\
Payload release & 4 metal-gear servos + mechanical detent & 1 & 4,500 & A power glitch must not drop a kit \\
Propellers & Tarot 1855 carbon, CW/CCW & 4 & 4,000 & See note below \\
Wiring & 10\,AWG silicone, XT90-S, JST-GH & 1 & 2,400 & Gauge is set by the 115\,A peak \\
Propeller adapters & 6\,mm self-tightening hubs & 4 & 1,400 & A nut loosening at 5{,}000\,rpm loses the aircraft \\
\bottomrule
\end{tabular}
\end{center}

\subsection*{Step 4 --- Fly aircraft one \quad \textbf{\rs{42{,}574}}}

\begin{center}\small
\begin{tabular}{@{}L{3.2cm}L{4.6cm}rrL{4.0cm}@{}}
\toprule
\textbf{Item} & \textbf{Model} & \textbf{Qty} & \textbf{Total} & \textbf{Why} \\
\midrule
Battery charger & \textbf{ToolkitRC M6D}, 500\,W / 25\,A, dual & 2 & 11,672 & See note below \\
Consumables & Fasteners, thread-lock, tape, filament & --- & 13,000 & Consumed continuously; running out stops work \\
Safety-pilot radio & \textbf{RadioMaster Boxer}, ExpressLRS & 1 & 6,955 & A pilot taking over manually needs full-size sticks \\
Sun hood & Fabricated, plus a second-hand monitor & 1 & 1,500 & Screens are unreadable in Punjab daylight \\
\bottomrule
\end{tabular}
\end{center}

\subsection*{Step 5 --- Ground station and registration \quad \textbf{\rs{66{,}992}}}

Brought forward ahead of the second and third aircraft, because centimetre
positioning has to be working before we can measure how accurately aircraft one
reports a person's coordinates.

\begin{center}\small
\begin{tabular}{@{}L{3.2cm}L{4.6cm}rrL{4.0cm}@{}}
\toprule
\textbf{Item} & \textbf{Model} & \textbf{Qty} & \textbf{Total} & \textbf{Why} \\
\midrule
Field consumables & Cables, mounts, field spares & --- & 25,000 & Corrections ride the existing radio, so no extra base radio \\
Base station & Second GNSS receiver on a tripod & 1 & 18,000 & Without it the aircraft have no corrections \\
Registration & Statutory, unmanned aircraft & 3 & 5,000 & At 6.36\,kg these are Small category under the Drone Rules 2021 \\
Survey tripod & Photographic tripod and adapter & 1 & 4,000 & Survey-grade is not required here \\
\bottomrule
\end{tabular}
\end{center}

\subsection*{Steps 6 and 7 --- Aircraft two and three \quad \textbf{\rs{1{,}99{,}539} each}}

Identical to steps 2 and 3 combined, repeated once per aircraft. These are the
two largest requests, and deliberately the last: by the time either is made, a
complete aircraft has already flown the mission, so what is being funded is
\emph{replication of a proven design} rather than a hope.

\subsection*{Two notes on the parts}

\textbf{A saving worth having.} The usual charger is the HOTA D6 Pro at
\rs{10{,}499}--\rs{13{,}500}, but at 325\,W per channel it can barely charge our
packs. The ToolkitRC M6D is 500\,W at \rs{5{,}836} --- cheaper and faster. It
needs a separate power supply, which the department already owns. Two units give
four channels, so the fleet charges at once on a flight day, and it still comes
in \rs{2{,}328} under budget.

\textbf{One line is under-budgeted, and we would rather say so now.} We carry
\rs{1{,}000} per propeller. The verified Indian listing is \rs{3{,}250} per
counter-rotating \emph{pair}, so four propellers is \rs{6{,}500} per aircraft
rather than \rs{4{,}000} --- about \rs{7{,}500} short across three aircraft, on
the item most often broken in testing. We will look for a cheaper supplier first
and absorb the rest from the charger saving and contingency; if neither works,
this line has to rise.

\subsection*{What splitting this way costs}

Smaller steps are not free. Buying one aircraft at a time means placing the same
imported order three times, and imported parts take four to six weeks. Ordering
all three sets together at step 2 would save roughly \textbf{eight to twelve
weeks} of schedule, at the price of a much larger single decision.

We have chosen the smaller steps because we would rather move slowly with your
confidence than quickly without it. If the schedule becomes the binding
constraint, consolidating steps 2, 6 and 7 into one order is the lever ---
and it is your call, not ours.

\textbf{For the detection test specifically}, we need nothing beyond what the
department already has: about \textbf{one week on a GPU} and \textbf{30\,GB of
storage}. The dataset is public and the software is open source.

\section{Conclusion}

The design is further along than it looks: the sizing, the error budgets and the
flight software are written down, reproducible, and honest about which numbers
are measured and which are still calculated.

What we are asking for is the means to finish it, released in four steps so that
nothing is committed before the previous step has produced evidence. The
immediate need is small --- about a week of GPU time to test the detection idea
on free public data.

We would rather find out now whether that idea holds than build on a figure we
calculated but never measured. Having already found one sizing mistake that way,
testing first seems clearly the better choice.

\vspace{4mm}
We would also value your guidance on four decisions that are ours to make but
that we would rather not make alone: the body size we design against, the
operating altitude, whether to adopt the two-stage detection approach, and which
of two processing modules to buy.

\end{document}
